% SOURCE_AUTHORITY: sga5_fr_workpass.tex, lines 12848--12868 (LF-normalized unit).
% SOURCE_SHA256: 791F4EFFC5E02832D5D77ED03518C8156D6F07E4C8238B03545DB93D883FBB28
\subsection*{6. Definição dos termos locais $\varepsilon_x^\Delta(F)$}

\subsubsection*{6.1}

Sejam $C$ uma curva algébrica conexa e lisa sobre o corpo algebricamente fechado $k$, $A$ uma $\Lambda$-álgebra ($\Lambda=\mathbb Z/\ell^\nu\mathbb Z$, com $\ell$ diferente da característica de $k$) e $\Delta\subset D^+(A)$ uma subcategoria triangulada de $D^+(A)$ estável sob somandos diretos. Seja $F$ um complexo de feixes de $A$-módulos sobre $C$ que satisfaça as seguintes condições:

i) Para todo ponto geométrico $\bar x\in C(k)$, tem-se $F_{\bar x}\in\Ob(\Delta)$.

ii) Existe um aberto não vazio de $C$ sobre o qual $F$ é localmente constante; ou, o que é equivalente, existem um aberto não vazio $U$ de $C$ e um recobrimento de Galois conexo $C'\xrightarrow{p}C$ de $C$, de grupo $G$, tais que, se $F'=p^*(F)$ é a imagem inversa de $F$ por $p$, então $F'|U'$ é constante $(U'=p^{-1}(U))$.

O feixe $F'|U'$ é definido, portanto, por um complexo $M$ de $A$-módulos sobre o qual opera $G$. Se se denota por $\eta$ (resp. $\eta'$) o ponto genérico de $C$ (resp. $C'$), e por $\bar\eta$ (resp. $\bar\eta'$) um ponto geométrico sobre $\eta$ (resp. $\eta'$), tem-se:
\[
p(\eta')=\eta,\qquad
F_{\bar\eta}\simeq F'_{\bar\eta'}=M.
\tag{6.1.1}
\]
Como complexo de $A$-módulos, $M$ é isomorfo a $F_{\bar\eta}$ e, graças à condição i), é um objeto de $\Delta$. Portanto, por (3.6), para todo $x\in C(k)$,
\[
\Hom^\bullet_{\Lambda[G]}(Sw_x^\Lambda,M)
\]
é um objeto de $\Delta$ e define, portanto, um elemento de $K(\Delta)$.
